\documentclass[12pt, a4paper]{ctexart}
\usepackage[margin=2cm]{geometry}
\usepackage{libertine}

\usepackage{titlesec}
\usepackage{zhnumber}
\titleformat*{\section}{\Large\bfseries\raggedright}
\renewcommand{\thesection}{\zhnum{section}}
\renewcommand{\thesubsection}{\arabic{subsection}}

\setcounter{secnumdepth}{4}
\renewcommand{\theparagraph}{\thesubsubsection.\arabic{paragraph}}
\renewcommand{\paragraph}[1]{
    \refstepcounter{paragraph}
    {\bfseries\theparagraph\quad#1}\par
    \vspace{2pt}
    \noindent
}

\usepackage{enumitem}
\setlist[enumerate]{itemsep=2pt, parsep=0pt, partopsep=0pt, topsep=2pt}
\setlist[itemize]{itemsep=2pt, parsep=0pt, partopsep=0pt, topsep=2pt}
\linespread{1.2}

\usepackage{graphicx}

\usepackage{minted}
\setminted{
    breaklines=true, escapeinside=||, fontsize=\small, frame=lines,
    mathescape=$$, numbers=none, style=bw, tabsize=4
}

\usepackage{siunitx}

\begin{document}
    
\pagestyle{plain}
\thispagestyle{empty}

\noindent
\begin{tabular*}{\textwidth}{l @{\extracolsep{\fill}} r @{\extracolsep{6pt}} l}
    \LARGE{\textbf{实验报告}} & 计算机网络 & \textit{Computer Networking} \\
\end{tabular*}\\\\
\begin{tabular*}{\textwidth}{l l}
    \textbf{报告标题: } & \textbf{BGP 路由协议验证} \\
    \textbf{学号: } & 19240212 \\
    \textbf{姓名: } & 华博文 \\
    \textbf{日期: } & 2025 年 12 月 29 日 \\
\end{tabular*}\\
\rule[2ex]{\textwidth}{2pt}

\section{实验目的}

本实验旨在掌握基于 IPMininet 的 BGP 路由协议配置方法，理解 BGP 作为域间动态路由协议的工作原理（包括自治系统（AS）标识、邻居建立、网络前缀通告等核心机制）；通过 IPMininet 搭建包含 2 台路由器、2 台主机的跨 AS 网络拓扑，完成设备 IP 地址配置、路由器 BGP 协议部署及跨 AS 网络连通性测试，强化对 BGP 在不同自治系统间路由自动化管理作用的认知，同时熟悉 IPMininet 仿真工具在域间路由实验中的应用流程。

\section{实验内容简要描述}

\begin{enumerate}
    \item \textbf{IPMininet 仿真网络启动}：在 Linux 环境中启动 Open vSwitch 服务，进入 IPMininet 示例目录，查看自定义 BGP 网络拓扑文件并启动包含 2 个 AS（AS 100、AS 200）的仿真网络；
    
    \item \textbf{主机与路由器 IP 配置}：配置主机网卡的 IP 地址及默认网关，配置路由器双网卡的 IP 地址以实现 AS 内部及 AS 间网段互联；
    
    \item \textbf{路由器 BGP 协议配置}：通过 \mintinline{bash}`telnet` 连接 BGP 进程，进入配置模式设置自治系统号、添加 BGP 邻居、通告本地 AS 的网络前缀，完成 BGP 协议的部署；
    
    \item \textbf{BGP 协议验证}：测试跨 AS 主机的 \mintinline{bash}`ping` 连通性，查看路由器路由表及 BGP 邻居 / 路由信息，验证 BGP 域间动态路由的生效状态。
\end{enumerate}

\section{实验步骤与结果分析}

\subsection{IPMininet 仿真网络启动}

右键打开 Linux 终端，执行命令启动 OVS 服务（为 IPMininet 提供交换功能支持）：

\begin{minted}{bash}
ovs-vswitchd --pidfile --detach
\end{minted}

进入 IPMininet 主目录（\mintinline{text}`/home/headless/ipmininet`），右键打开终端，执行命令启动目标拓扑的仿真网络：

\begin{minted}{bash}
sudo python3 -m ipmininet.examples --topo=simple_bgp4_network
\end{minted}

终端依次显示设备添加日志，最终进入 Mininet CLI 界面，确认网络启动完成。

\begin{figure}[H]
    \centering
    \includegraphics[width=\textwidth]{./img/screenshot_0000.png}
\end{figure}

\subsection{主机与路由器 IP 配置}

在 Mininet CLI 中打开 h1 终端：

\begin{minted}{bash}
xterm h1
\end{minted}

执行 \mintinline{bash}`ifconfig` 查看默认 IP（与实验规划不符），执行命令将 h1 的网卡 IP 配置为 AS 100 网段（\mintinline{text}`10.0.0.2`）：

\begin{minted}{bash}
sudo ifconfig h1-eth0 10.0.0.2
\end{minted}

在 Mininet CLI 中打开 h2 终端：

\begin{minted}{bash}
xterm h2
\end{minted}

同理，执行命令将 h2 的网卡 IP 配置为 AS 200 网段（\mintinline{text}`192.168.0.2`）：

\begin{minted}{bash}
sudo ifconfig h2-eth0 192.168.0.2
\end{minted}

\begin{figure}[H]
    \centering
    \includegraphics[width=\textwidth]{./img/screenshot_0001.png}
\end{figure}

在 Mininet CLI 中打开路由器 r1（属于 AS 100）终端：

\begin{minted}{bash}
xterm r1
\end{minted}

执行 \mintinline{bash}`ifconfig` 查看网卡信息，配置双网卡 IP（\mintinline{text}`10.0.0.1/24` 为 AS 100 内部网段，\mintinline{text}`10.0.5.1/24` 为 AS 间互联网段）：

\begin{minted}{bash}
sudo ifconfig r1-eth0 10.0.0.1/24
sudo ifconfig r1-eth1 10.0.5.1/24
\end{minted}

在 Mininet CLI 中打开路由器 r2（属于 AS 200）终端，配置双网卡 IP：

\begin{minted}{bash}
sudo ifconfig r2-eth0 10.0.5.2/24
sudo ifconfig r2-eth1 192.168.0.1/24
\end{minted}

\begin{figure}[H]
    \centering
    \includegraphics[width=\textwidth]{./img/screenshot_0002.png}
\end{figure}

在 h1 终端中添加默认网关（指向 r1 的 AS 内部接口）：

\begin{minted}{bash}
sudo route add default gw 10.0.0.1
\end{minted}

在 h2 终端中添加默认网关（指向 r2 的 AS 内部接口）：

\begin{minted}{bash}
sudo route add default gw 192.168.0.1
\end{minted}

执行 \mintinline{bash}`route -n` 确认配置生效：

\begin{figure}[H]
    \centering
    \includegraphics[width=\textwidth]{./img/screenshot_0003.png}
\end{figure}

\subsection{配置路由器 BGP 协议}

在 Mininet CLI 中查看 r1 路由表，发现无跨 AS 网段路由：

\begin{minted}{bash}
r1 route
\end{minted}

\begin{figure}[H]
    \centering
    \includegraphics[width=\textwidth]{./img/screenshot_0004.png}
\end{figure}

在 r1 终端中通过 \mintinline{bash}`telnet` 连接 BGP 进程（默认密码：\mintinline{text}`zebra`）：

\begin{minted}{bash}
telnet localhost bgpd
\end{minted}

进入 BGP 配置模式，设置 AS 号为 100，并添加 BGP 邻居（r2 的 AS 间接口 IP \mintinline{text}`10.0.5.2`，对应 AS 200）：

\begin{minted}{bash}
enable
configure terminal
router bgp 100
no bgp ebgp-requires-policy
neighbor 10.0.5.2 remote-as 200
neighbor 10.0.5.2 description "router 2"
write
\end{minted}

执行 \mintinline{bash}`show ip bgp summary` 查看 BGP 邻居状态（初始为 active，表示正在建立邻居）。

在 r2 终端中重复上述步骤，配置 BGP（AS 号为 200，邻居为 r1 的 AS 间接口 IP \mintinline{text}`10.0.5.1`，对应 AS 100）：

\begin{minted}{bash}
enable
configure terminal
router bgp 200
no bgp ebgp-requires-policy
neighbor 10.0.5.1 remote-as 100
neighbor 10.0.5.1 description "router 1"
write
\end{minted}

查看 r2 的 BGP 邻居中间状态：

\begin{figure}[H]
    \centering
    \includegraphics[width=\textwidth]{./img/screenshot_0005.png}
\end{figure}

从结果可见，BGP 邻居处于 active 或 connect 状态（正在建立连接），BGP 路由表为空，无任何网络前缀信息，这是未通告本地网段时的中间状态。

返回 r1 的 BGP 终端，重新进入配置模式，添加本地 AS 网段通告并保存配置：

\begin{minted}{bash}
configure terminal
router bgp 100
network 10.0.0.0/24
write
\end{minted}

再次查看 r1 的 BGP 状态，确认网段已通告及邻居状态变化：

\begin{minted}{bash}
show ip bgp summary
\end{minted}

在 r2 终端中重新进入 BGP 配置模式，补充通告 AS 200 本地网段并保存配置：

\begin{minted}{bash}
configure terminal
router bgp 200
network 192.168.0.0/24
write
\end{minted}

查看 r2 的 BGP 邻居状态，显示 Up 表示邻居建立成功，且已识别本地通告网段：

\begin{minted}{bash}
show ip bgp summary
\end{minted}

\begin{figure}[H]
    \centering
    \includegraphics[width=\textwidth]{./img/screenshot_0006.png}
\end{figure}

\subsection{验证 BGP 协议}

在 h1 终端中 ping h2 的 IP（跨 AS 通信），结果显示连通（无丢包）：

\begin{minted}{bash}
ping 192.168.0.2
\end{minted}

在 Mininet CLI 中查看 r1、r2 的路由表，可见 BGP 学习到的跨 AS 动态路由（标志 UG 表示网络可达且邻居为网关）：

\begin{minted}{bash}
r1 route
r2 route
\end{minted}

\begin{figure}[H]
    \centering
    \includegraphics[width=\textwidth]{./img/screenshot_0007.png}
\end{figure}

\section{实验中遇到的问题及体会}

在启动 IPMininet 时，因未提前启动 OVS 服务，终端提示 “交换设备不可用” 错误。执行 \mintinline{text}`ovs-vswitchd --pidfile --detach` 启动服务后，网络成功启动 —— 这让我认识到 IPMininet 依赖 OVS 提供交换能力，环境准备是实验的前提。

配置 BGP 邻居时，最初遗漏 \mintinline{text}`no bgp ebgp-requires-policy` 命令，导致邻居始终无法建立。查阅资料后得知，该命令用于关闭 BGP 的策略检查限制（实验环境下需临时关闭），执行后邻居状态转为 Up —— 这体现了 BGP 协议的安全性设计，同时也说明实验配置需适配工具的默认规则。

本次实验分两步配置 BGP 协议（先配邻居看中间状态，再通告网段），让我清晰观察到 BGP 配置的阶段性效果：未通告网段时，邻居仅处于连接建立阶段，路由表为空，无法实现跨 AS 通信；补充网段通告后，BGP 才会将本地网段传递给邻居，进而生成跨 AS 路由。这种分步操作的方式，帮助我更直观地理解了 BGP 邻居建立与网段通告的独立关联性，而非一次性配置的黑盒操作。

通过本次实验，我直观体会到 BGP 作为域间路由协议的核心作用：不同 AS 间无需手动配置静态路由，仅需建立 BGP 邻居并通告本地网段，即可自动学习跨 AS 路由。相比 OSPF（域内路由），BGP 更注重 “自治系统间的路由协商”，这与 AS 的管理边界特性相匹配。

实验也深化了 “跨域路由协同” 的认知：从 AS 号的区分、AS 间互联网段的规划，到 BGP 邻居的建立与网段通告，每一步都需严格匹配双方配置，任何一处参数（如 AS 号、邻居 IP）错误都会导致跨域通信失效。这种 “域间协同” 的逻辑，是互联网跨自治系统通信的基础。

最后，IPMininet 的虚拟环境为 BGP 实验提供了低成本的实践载体，可快速搭建跨 AS 拓扑并反复调试配置。这一过程不仅提升了我的 Linux 命令与 BGP 配置熟练度，更培养了从路由表、BGP 邻居状态倒推故障的排查思维，为后续复杂域间路由技术的学习奠定了基础。

\end{document}